% !TeX program = pdflatex
\documentclass[aspectratio=169,xcolor={dvipsnames,table}]{beamer}

\input{../../_en_preambulo_beamer}
% Comandos y macros estadísticas compartidas para las presentaciones Beamer en inglés (Metropolis)

% Conjuntos numéricos básicos
\newcommand{\R}{\mathbb{R}}
\newcommand{\N}{\mathbb{N}}
\newcommand{\Q}{\mathbb{Q}}
\newcommand{\Z}{\mathbb{Z}}
\newcommand{\set}[1]{\left\{ #1 \right\}}
\newcommand{\abs}[1]{\left|#1\right|}
\newcommand{\norm}[1]{\left\| #1 \right\|}

% Comandos estadísticos y probabilísticos del curso
\newcommand{\Var}{\operatorname{Var}}
\newcommand{\cov}{\operatorname{Cov}}
\newcommand{\corr}{\rho}
\newcommand{\comb}[2]{\begin{pmatrix} #1 \\ #2 \end{pmatrix}}
\newcommand{\s}{\sigma}
\newcommand{\particion}{\mathfrak{P}}
\newcommand{\card}[1]{\#\left( #1 \right)}
\newcommand{\seq}[3]{\set{#1}_{#2}^{#3}}
\newcommand{\E}{\mathbb{E}}
\newcommand{\Prob}{\mathbb{P}}

% Entornos pedagógicos y cajas en inglés académico natural (soportando tanto nombres en inglés como en español)
\newenvironment{discussion}{%
  \begin{alertblock}{Discussion Question}%
}{%
  \end{alertblock}%
}
\newenvironment{discusion}{%
  \begin{alertblock}{Discussion Question}%
}{%
  \end{alertblock}%
}

\newenvironment{reflection}{%
  \begin{alertblock}{Classroom Reflection}%
}{%
  \end{alertblock}%
}
\newenvironment{reflexion}{%
  \begin{alertblock}{Classroom Reflection}%
}{%
  \end{alertblock}%
}

\newenvironment{paradox}{%
  \begin{alertblock}{Exploratory Paradox}%
}{%
  \end{alertblock}%
}
\newenvironment{paradoja}{%
  \begin{alertblock}{Exploratory Paradox}%
}{%
  \end{alertblock}%
}

\newenvironment{motivation}{%
  \begin{exampleblock}{Motivation: Data Science in Practice}%
}{%
  \end{exampleblock}%
}
\newenvironment{motivacion}{%
  \begin{exampleblock}{Motivation: Data Science in Practice}%
}{%
  \end{exampleblock}%
}

\newenvironment{quickchallenge}{%
  \begin{block}{Quick Team Challenge}%
}{%
  \end{block}%
}
\newenvironment{retoclase}{%
  \begin{block}{Quick Team Challenge}%
}{%
  \end{block}%
}


\title{Discrete Distributions in Data Science}
\subtitle{Section 03.10 --- MLE Fitting, Overdispersion Detection, and Information Criteria}
\author[J. Castillo Colmenares]{Juliho Castillo Colmenares}
\institute[Tec de Monterrey]{Tecnologico de Monterrey}
\date{\vspace{-1.5cm}}

\begin{document}

% SLIDE 01: Title Page
\begin{frame}[plain]
	\titlepage
\end{frame}

% SLIDE 02: Roadmap
\begin{frame}{Roadmap --- Chapter 03: Discrete Random Variables}
	\vspace{-0.15cm}
	\begin{block}{Curricular Structure of Chapter 03}
		\footnotesize
		\begin{enumerate}\itemsep=0.01cm
			\item \textcolor{gray}{Section 03.01: Discrete Probability Mass Functions (PMF and Support)}
			\item \textcolor{gray}{Section 03.02: Cumulative Distribution Functions (Discrete CDF)}
			\item \textcolor{gray}{Section 03.03: Mathematical Expectation, Variance, and Moments}
			\item \textcolor{gray}{Section 03.04: Bernoulli and Binomial Distributions}
			\item \textcolor{gray}{Section 03.05: Geometric and Negative Binomial Distributions}
			\item \textcolor{gray}{Section 03.06: Hypergeometric Distribution}
			\item \textcolor{gray}{Section 03.07: Poisson Distribution}
			\item \textcolor{gray}{Section 03.08: Multinomial Distribution}
			\item \textcolor{gray}{Section 03.09: Normal Distribution}
			\item \textbf{\textcolor{TecRojo}{Section 03.10: Discrete Distributions in Data Science}} $\leftarrow$ \emph{Current focus and Chapter closing}
		\end{enumerate}
	\end{block}
	\vspace{-0.3cm}
	\begin{alertblock}{Learning Objective}
		\scriptsize
		Integrate all discrete distributions into modern applications: MLE fitting, overdispersion detection, information criteria for model selection, and likelihood ratio inference.
	\end{alertblock}
\end{frame}

% SLIDE 02B: Motivation
\begin{frame}{From Theory to Practice: Distributions in Real Data}
	\vspace{-0.15cm}
	\begin{block}{Why Fit Distributions to Empirical Data?}
		\footnotesize
		In practice, we rarely know with certainty which distribution generated our data. What we observe are counts (clicks, defects, cancellations, calls), and we need to infer the model that best describes them. \emph{Maximum likelihood} (MLE) is the standard method: given a parametric family, we choose the parameters that maximize the probability of observing the data.
	\end{block}
	\vspace{-0.1cm}
	\pause
	\begin{alertblock}{Key Ideas}
		\footnotesize
		\begin{itemize}\itemsep=0.06cm
			\item \textbf{Parametric fitting:} Poisson assumes equidispersion; if the empirical variance is larger, Negative Binomial is preferable.
			\item \textbf{Model selection:} AIC, BIC, and the likelihood ratio test balance goodness of fit and parsimony.
			\item \textbf{Uncertainty:} Confidence intervals and bootstrap quantify the variability of point estimators.
		\end{itemize}
	\end{alertblock}
\end{frame}

% SLIDE 03: MLE and Overdispersion
\begin{frame}{MLE Fitting and Overdispersion Detection}
	\vspace{-0.15cm}
	\begin{columns}[T]
		\begin{column}{0.48\textwidth}
			\begin{block}{Poisson MLE}
				\footnotesize
				For count data $x_1, \dots, x_n$ with $X_i \sim \text{Pois}(\lambda)$:
				$$L(\lambda) = \prod_{i=1}^{n} \dfrac{e^{-\lambda}\lambda^{x_i}}{x_i!}$$
				$\ell'(\lambda) = -n + \sum x_i / \lambda = 0 \implies \hat{\lambda} = \bar{x}$.
			\end{block}
		\end{column}
		\begin{column}{0.48\textwidth}
			\begin{alertblock}{Dispersion Index $D = s^2/\bar{x}$}
				\footnotesize
				\begin{itemize}\itemsep=0.04cm
					\item $D \approx 1$: equidispersion; Poisson adequate.
					\item $D > 1.5$: strong overdispersion; use Negative Binomial.
					\item $D < 0.8$: underdispersion; use Conway-Maxwell-Poisson.
				\end{itemize}
			\end{alertblock}
		\end{column}
	\end{columns}
\end{frame}

% SLIDE 03B: Negative Binomial as Alternative
\begin{frame}{Negative Binomial: The Natural Alternative for Overdispersion}
	\vspace{-0.15cm}
	\begin{block}{Negative Binomial PMF}
		\footnotesize
		The Negative Binomial distribution introduces an additional parameter $r > 0$ to model extra variability:
		$$P(X = k) = \binom{k + r - 1}{k} p^{r} (1-p)^{k}, \quad k = 0, 1, 2, \dots$$
		Mean: $\E[X] = r(1-p)/p = \mu$. Variance: $\Var(X) = \mu + \mu^{2}/r$, capturing overdispersion.
	\end{block}
	\vspace{-0.1cm}
	\pause
	\begin{alertblock}{Hierarchical Interpretation}
		\footnotesize
		\begin{itemize}\itemsep=0.03cm
			\item Under the model $\lambda_{i} \sim \text{Gamma}(r, p)$ with $X_{i} \mid \lambda_{i} \sim \text{Poisson}(\lambda_{i})$, marginally $X_{i} \sim \text{NegBin}(r, p)$.
			\item Method-of-moments MLE: $\hat{r} = \bar{x}^{2}/(s^{2} - \bar{x})$ and $\hat{p} = \bar{x}/s^{2}$.
			\item As $r \to \infty$, the Negative Binomial collapses to the Poisson (Problem 3.10.7).
		\end{itemize}
	\end{alertblock}
\end{frame}

% SLIDE 04: Information Criteria
\begin{frame}{Information Criteria: AIC and BIC for Model Selection}
	\vspace{-0.15cm}
	\begin{columns}[T]
		\begin{column}{0.48\textwidth}
			\begin{block}{Akaike and Bayesian Criteria}
				\footnotesize
				$$\text{AIC} = 2k - 2\ell(\hat{\theta}), \quad \text{BIC} = k \log n - 2\ell(\hat{\theta})$$
				Choose model with lower AIC/BIC. Differences $\Delta < 2$ are insignificant; $\Delta > 10$ are very strong.
			\end{block}
		\end{column}
		\begin{column}{0.48\textwidth}
			\begin{alertblock}{Wilks Likelihood Ratio Test}
				\footnotesize
				For nested models $L_0 \subset L_1$ with $k_0 < k_1$ parameters:
				$$\Lambda = -2(\log L_0(\hat{\theta}_0) - \log L_1(\hat{\theta}_1)) \xrightarrow{d} \chi^2_{k_1 - k_0}$$
				Reject $H_0$ (restricted model) if $\Lambda > \chi^2_{k_1 - k_0, 0.05}$.
			\end{alertblock}
		\end{column}
	\end{columns}
\end{frame}

% SLIDE 05: Confidence Intervals and Bootstrap
\begin{frame}{Confidence Intervals and Bootstrap for Discrete Parameters}
	\vspace{-0.15cm}
	\begin{columns}[T]
		\begin{column}{0.48\textwidth}
			\begin{block}{Wald and Exact (Garwood) Intervals}
				\footnotesize
				Wald: $\hat{\lambda} \pm z_{0.025}\sqrt{\hat{\lambda}/n}$ (valid for $\hat{\lambda} \ge 30$). Exact: $\frac{1}{2n}\chi^2_{0.025, 2k}$ and $\frac{1}{2n}\chi^2_{0.975, 2(k+1)}$.
			\end{block}
		\end{column}
		\begin{column}{0.48\textwidth}
			\begin{alertblock}{Non-parametric Bootstrap}
				\footnotesize
				Resample with replacement $B$ times, compute $B$ MLEs, and use $\text{SE}_{\text{boot}} = \text{std}(\hat{\theta}^{*}_b)$. Percentile CI: $[\hat{\theta}^{*}_{0.025}, \hat{\theta}^{*}_{0.975}]$.
			\end{alertblock}
		\end{column}
	\end{columns}
\end{frame}

% SLIDE 06: Applications
\begin{frame}{Applications in Data Science and Business Analytics}
	\vspace{-0.1cm}
	\begin{itemize}\itemsep=0.1cm
		\item \textbf{Churn Analysis:} Modeling customer cancellations with Poisson and Negative Binomial to identify risk factors. \pause
		\item \textbf{Fraud Detection:} Counting suspicious transactions; overdispersion flags anomalous behavior. \pause
		\item \textbf{A/B Experimentation:} Comparing conversion rates (Binomial) and event counts (Poisson/NegBin) between variants. \pause
		\item \textbf{Industrial Reliability:} Time-to-failure (Exponential/Weibull) and defects per unit (Poisson/NegBin).
	\end{itemize}
\end{frame}

% SLIDE 07: Python Lab Block 1
\begin{frame}[fragile]{Python Lab: MLE and Overdispersion}
	\vspace{-0.05cm}
	\scriptsize
	Poisson and Negative Binomial MLE fitting with AIC/BIC comparison (\texttt{03.10\_discrete\_distributions\_data\_science.py}):
	\vspace{0.02cm}
	{\lstset{basicstyle=\fontsize{4.5pt}{5.5pt}\selectfont\ttfamily}
	\lstinputlisting[firstline=15, lastline=40]{../../code/03_variables_aleatorias_discretas/03.10_discrete_distributions_data_science.py}}
\end{frame}

% SLIDE 08: Python Lab Block 2
\begin{frame}[fragile]{Python Lab: Wilks Test and Overdispersion}
	\vspace{-0.05cm}
	\scriptsize
	Likelihood ratio test for detecting significant overdispersion:
	\vspace{0.02cm}
	{\lstset{basicstyle=\fontsize{4.5pt}{5.5pt}\selectfont\ttfamily}
	\lstinputlisting[firstline=48, lastline=72]{../../code/03_variables_aleatorias_discretas/03.10_discrete_distributions_data_science.py}}
\end{frame}

% SLIDE 09: Python Lab Block 3
\begin{frame}[fragile]{Python Lab: CIs and Bootstrap}
	\vspace{-0.05cm}
	\scriptsize
	Three CI methods (Wald, exact, bootstrap) and bootstrap for Negative Binomial parameters:
	\vspace{0.02cm}
	{\lstset{basicstyle=\fontsize{4.5pt}{5.5pt}\selectfont\ttfamily}
	\lstinputlisting[firstline=81, lastline=108]{../../code/03_variables_aleatorias_discretas/03.10_discrete_distributions_data_science.py}}
\end{frame}

% SLIDE 10: Python Lab Terminal Output
\begin{frame}[fragile]{Python Lab: Terminal Output}
	\vspace{-0.1cm}
	\begin{block}{}
		\fontsize{4pt}{5pt}\selectfont
		\begin{verbatim}
=== Block 1: MLE Fitting & Overdispersion ===
NegBin(r=2, p=0.4) | mean=2.95 | var=5.24 | D=1.78
  -> Strong overdispersion (D > 1.5): NegBin recommended
Poisson: logLik=-220.16, AIC=442.31, BIC=444.92
NegBin:  logLik=-211.05, AIC=426.10, BIC=431.31
  Delta AIC (P-NB): 16.21 | Delta BIC: 13.60

=== Block 2: Likelihood Ratio Test & Wilks ===
NegBin(r=3, p=0.5) n=200: mean=2.83, var=5.29
  Poisson logLik=-444.31 | NegBin logLik=-421.64
  Lambda = 45.34 | p-value = 0.000000
  Decision: Reject H0. NegBin significantly better.
LR vs overdispersion (n=100):
  r=10.0 (v/m=1.33): LR=4.24, p=0.040
  r=2.0  (v/m=1.58): LR=12.81, p=0.000
  r=1.0  (v/m=2.85): LR=78.66, p=0.000
  r=0.5  (v/m=7.99): LR=310.80, p=0.000

=== Block 3: Confidence Intervals & Bootstrap ===
12 months claims: sample mean=46.50
  Wald 95% CI: [42.64, 50.36]
  Exact 95% CI: [42.72, 50.52]
  Bootstrap SE: 1.97 | Wald SE: 1.97
  Bootstrap 95% CI: [42.75, 50.42]
NegBin bootstrap (n=200 churn data):
  Method of moments: r=2.03, p=0.40
  Bootstrap mean: r=2.13 +- 0.42 | p=0.41 +- 0.05
		\end{verbatim}
	\end{block}
\end{frame}

% SLIDE 11: Exercise Level 1 (Statement)
\begin{frame}{In-Class Exercise --- Level 1: Fundamental (1/2)}
	\vspace{-0.1cm}
	\begin{block}{Problem 3.10.1 --- Overdispersion Detection in Ads (Statement)}
		A web analytics team monitors the number of clicks received by $n = 100$ ads in one hour. Observed data have $\bar{x} = 4.2$ and $s^{2} = 8.7$. Calculate the Poisson MLE $\hat{\lambda}$ and the dispersion index $D = s^2/\bar{x}$; interpret whether the data exhibit overdispersion, underdispersion, or equidispersion.
	\end{block}
	\vspace{0.15cm}
	\pause
	\begin{alertblock}{Model Setup}
		\begin{itemize}\itemsep=0.04cm
			\item Log-likelihood: $\ell(\lambda) = -n\lambda + (\sum x_i)\log \lambda - \sum \log x_i!$.
			\item Derivative: $\ell'(\lambda) = -n + \sum x_i / \lambda = 0 \implies \hat{\lambda} = \bar{x} = 4.2$.
			\item Diagnostic: $D = s^2/\bar{x} = 8.7/4.2 \approx 2.07$.
		\end{itemize}
	\end{alertblock}
\end{frame}

% SLIDE 12: Exercise Level 1 (Solution)
\begin{frame}{In-Class Exercise --- Level 1: Fundamental (2/2)}
	\vspace{-0.05cm}
	\scriptsize
	The Poisson MLE is the sample mean, and the dispersion index reveals overdispersion: \pause
	\begin{block}{Point Estimation}
		\begin{align*}
			\hat{\lambda} = \bar{x} = 4.2 \quad \text{(Poisson MLE)}
		\end{align*}
	\end{block}
	\vspace{0.05cm}
	\pause
	\begin{alertblock}{Dispersion Index and Interpretation}
		\scriptsize
		$D = s^2/\bar{x} = 8.7/4.2 \approx 2.07 \gg 1$. The data exhibit \textbf{significant overdispersion}: the empirical variance is more than twice the mean. Under Poisson, we would expect $D \approx 1$. The pure Poisson model is structurally inadequate. Alternatives: Negative Binomial, robust-variance Poisson regression, or hierarchical model with random effects.
	\end{alertblock}
\end{frame}

% SLIDE 13: Exercise Level 2 (Statement)
\begin{frame}{In-Class Exercise --- Level 2: Operational (1/2)}
	\vspace{-0.1cm}
	\begin{block}{Problem 3.10.5 --- Poisson vs. Negative Binomial Comparison (Statement)}
		A streaming platform analyzes daily views per user. From 200 observations: $\bar{x} = 8.5$ and $s^{2} = 9.8$. Fit Poisson and Negative Binomial models, compute AIC and BIC, and decide which is preferable.
	\end{block}
	\vspace{0.15cm}
	\pause
	\begin{alertblock}{Comparison Metrics}
		\scriptsize
		Use AIC $= 2k - 2\ell(\hat{\theta})$ and BIC $= k\log n - 2\ell(\hat{\theta})$. The model with lower AIC (or BIC) is preferable by parsimony. Differences $\Delta\text{AIC} < 2$ are insignificant; $\Delta\text{AIC} > 10$ are very strong.
	\end{alertblock}
\end{frame}

% SLIDE 14: Exercise Level 2 (Solution)
\begin{frame}{In-Class Exercise --- Level 2: Operational (2/2)}
	\scriptsize
	Compute the MLEs and information criteria: \pause
	\begin{block}{Model MLEs}
		\scriptsize
		Poisson: $\hat{\lambda} = 8.5$. Negative Binomial (method of moments):
		$$\hat{r} = \dfrac{\bar{x}^{2}}{s^{2} - \bar{x}} = \dfrac{72.25}{1.3} \approx 55.58, \quad \hat{p} = \dfrac{\bar{x}}{s^{2}} \approx 0.867$$
	\end{block}
	\vspace{0.05cm}
	\pause
	\begin{alertblock}{Parsimony Decision}
		\scriptsize
		With $\hat{r} = 55.58 \gg 1$, the Negative Binomial collapses practically to Poisson (Problem 3.10.7). The dispersion index $s^2/\bar{x} = 1.15$ indicates mild overdispersion. The expected AIC difference is $< 2$, suggesting preference for Poisson by parsimony. \emph{Conclusion:} with mild overdispersion, the simpler model (Poisson) is preferable.
	\end{alertblock}
\end{frame}

% SLIDE 15: Exercise Level 3 (Statement)
\begin{frame}{In-Class Exercise --- Level 3: Analytical (1/2)}
	\vspace{-0.1cm}
	\begin{block}{Problem 3.10.7 --- Negative Binomial Collapse to Poisson (Statement)}
		Prove analytically that the Negative Binomial $\text{NB}(r, p)$ collapses to $\text{Pois}(\lambda)$ as $r \to \infty$ with $\lambda = r(1-p)/p$ fixed, and discuss implications for model selection.
	\end{block}
	\vspace{0.1cm}
	\pause
	\begin{alertblock}{Deduction Strategy}
		\scriptsize
		Start with the NB PMF, substitute $p = r/(r + \lambda)$ and take the limit $r \to \infty$ with $\lambda$ fixed. Use that for large $r$ with $j$ fixed, $r(r+1)\cdots(r+j-1) \approx r^{j}$ and $(1 - \lambda/r)^{r} \to e^{-\lambda}$.
	\end{alertblock}
\end{frame}

% SLIDE 16: Exercise Level 3 (Solution)
\begin{frame}{In-Class Exercise --- Level 3: Analytical (2/2)}
	\vspace{-0.05cm}
	\scriptsize
	Take the limit $r \to \infty$ with $\lambda$ fixed: \pause
	\begin{block}{Asymptotic Expansion}
		\scriptsize
		With $p = r/(r + \lambda) \approx 1 - \lambda/r$ and $1 - p = \lambda/(r + \lambda) \approx \lambda/r$ as $r \to \infty$:
		\begin{align*}
			P(X = k) &= \binom{k + r - 1}{k} p^{r} (1-p)^{k} \\
			&\approx \dfrac{r^{k}}{k!} \cdot \left(1 - \dfrac{\lambda}{r}\right)^{r} \cdot \left(\dfrac{\lambda}{r}\right)^{k} \to \dfrac{1}{k!} \cdot e^{-\lambda} \cdot \lambda^{k} = \dfrac{e^{-\lambda}\lambda^{k}}{k!}
		\end{align*}
	\end{block}
	\vspace{0.05cm}
	\pause
	\begin{alertblock}{Model Selection Implication}
		\scriptsize
		If the MLE $\hat{r}$ is very large, the data are consistent with Poisson and the Negative Binomial adds little value. AIC/BIC will correctly penalize the extra parameter. Practical rule: if $\hat{r} > 100$, use Poisson for parsimony. \qedhere
	\end{alertblock}
\end{frame}

% SLIDE 17: Exercise Level 4 (Statement)
\begin{frame}{In-Class Exercise --- Level 4: Challenging (1/2)}
	\vspace{-0.1cm}
	\begin{block}{Problem 3.10.10 --- Hierarchical Poisson-Gamma Model (Statement)}
		An insurance firm observes claim counts with $\bar{x} = 1.2$ and $s^{2} = 3.4$. Management suspects unobserved heterogeneity: each policyholder has rate $\lambda_{i} \sim \text{Gamma}(\alpha, \beta)$. Prove that marginally $X_{i} \sim \text{NB}(r, p)$ and estimate parameters by the method of moments.
	\end{block}
	\vspace{0.15cm}
	\pause
	\begin{alertblock}{Application}
		\scriptsize
		Compute the likelihood ratio statistic to compare with Poisson and conclude whether unobserved heterogeneity is statistically significant.
	\end{alertblock}
\end{frame}

% SLIDE 18: Exercise Level 4 (Solution)
\begin{frame}{In-Class Exercise --- Level 4: Challenging (2/2)}
	\vspace{-0.05cm}
	\scriptsize
	Integrate over the Poisson-Gamma mixture to obtain the Negative Binomial: \pause
	\begin{block}{Marginal Derivation}
		\scriptsize
		$P(X = k) = \int_{0}^{\infty} \dfrac{e^{-\lambda}\lambda^{k}}{k!} \cdot \dfrac{\beta^{\alpha}}{\Gamma(\alpha)}\lambda^{\alpha-1}e^{-\beta\lambda}\,d\lambda = \binom{k + \alpha - 1}{k}\left(\dfrac{1}{\beta+1}\right)^{\alpha}\left(\dfrac{\beta}{\beta+1}\right)^{k}$
	\end{block}
	\vspace{0.05cm}
	\pause
	\begin{alertblock}{Method of Moments Estimation}
		\scriptsize
		With $r = \alpha$ and $p = 1/(\beta+1)$:
		\begin{align*}
			\hat{r} = \dfrac{\bar{x}^{2}}{s^{2} - \bar{x}} = \dfrac{1.44}{2.2} \approx 0.655, \quad \hat{p} = \dfrac{\bar{x}}{s^{2}} = \dfrac{1.2}{3.4} \approx 0.353.
		\end{align*}
		With $s^2/\bar{x} = 3.4/1.2 \approx 2.83 \gg 1$, the Wilks test gives $\Lambda \gg 3.841$ ($p < 0.001$). \emph{Conclusion:} unobserved heterogeneity is statistically significant; Negative Binomial is preferable to Poisson. \qedhere
	\end{alertblock}
\end{frame}

% SLIDE 19: Closing Chapter 03
\begin{frame}{Closing Chapter 03 and Modular Perspective}
	\vspace{-0.3cm}
	\begin{block}{Synthesis: Discrete Distributions Toolkit}
		\scriptsize
		Over 10 sections, we mastered the ecosystem of discrete distributions: Binomial, Geometric, Hypergeometric, Poisson, Multinomial, and their Normal approximation. We close with inference tools: MLE, information criteria, likelihood ratio tests, and bootstrap.
	\end{block}
	\vspace{0.0cm}
	\pause
	\begin{alertblock}{Key Lessons and Next Step}
		\begin{itemize}\itemsep=0.02cm
			\item \textbf{Complete Family:} Binomial, Geometric, Hypergeometric, Poisson, Multinomial, Normal.
			\item \textbf{Equidispersion vs. Overdispersion:} Poisson assumes $\E[X] = \Var(X)$; NegBin relaxes this.
			\item \textbf{Rigorous Inference:} AIC/BIC balance goodness of fit and parsimony; Wilks validates nested models.
			\item \textbf{Next Chapter:} \textcolor{TecRojo}{Continuous Random Variables}.
		\end{itemize}
	\end{alertblock}
\end{frame}

\end{document}
